% Source: Evans, "Partial Differential Equations" (2nd ed., AMS Graduate Studies in Mathematics vol. 19, 2010),
% Chapter E "Measure Theory", PDF pages 647-652 of MathBooks/Evans Partial Differential Equations(1).pdf.
% Transcribed page-by-page by vision subagents, then merged and restructured into
% leanblueprint conventions (semantic \label, bracketed titles, \uses dependency edges) below.
%
% NOTE: PDF page 647 opens with the tail end of Appendix D "Linear Functional Analysis"
% -- Theorem 7 on the existence of an orthonormal eigenbasis for a compact symmetric
% operator, continuing directly from appendixD.tex's final result (the lemma on bounds
% on the spectrum of a symmetric operator). This leftover fragment is transcribed here,
% at the point where it physically occurs in the source page range handed to this pass,
% immediately before the "APPENDIX E: MEASURE THEORY" title. It uses the \dcref{appD:...}
% numbering of Appendix D (it is literally Appendix D's Theorem 7), not \dcref{appE:...}.
%
% NOTE ON GAPS (resolved): PDF page 650 originally returned a vision-subagent refusal
% placeholder rather than actual page content, leaving a gap between the "integrable"/
% "summable" definitions (end of \S E.2) and Theorem 6 (Lebesgue's Differentiation
% Theorem, start of \S E.4). That gap has since been closed by directly re-reading the
% source page image: it contained the essential-supremum notation closing \S E.2,
% then the entirety of \S E.3 "Convergence Theorems for Integrals" (Fatou's Lemma, the
% Monotone Convergence Theorem, and the Dominated Convergence Theorem), then the opening
% sentence of \S E.4 "Differentiation" (there is no separate Fubini/product-measures
% section -- that was only a placeholder guess, not the book's actual content).

\begin{theorem}[Eigenvectors of a compact, symmetric operator]
\label{thm:eigenvectors-compact-symmetric-operator}
\dcref{appD:thm-7}
\group{Fredholm Theory}
\uses{def:hilbert-space,def:separable-space,def:compact-linear-operator,def:adjoint-symmetric-operator,def:eigenvalue-eigenvector-appendix-d,thm:fredholm-alternative,thm:spectrum-of-compact-operator,lem:bounds-on-spectrum-symmetric-operator}
\textit{Let $H$ be a separable Hilbert space, and suppose $S : H \to H$ is a compact and symmetric operator. Then there exists a countable orthonormal basis of $H$ consisting of eigenvectors of $S$.}
\end{theorem}

\begin{proof}
1. Let $\{\eta_k\}$ comprise the sequence of distinct eigenvalues of $S$, excepting 0. Set $\eta_0 = 0$. Write $H_0 = N(S)$, $H_k = N(S - \eta_k I)$ ($k = 1, \ldots$). Then $0 \leq \dim H_0 \leq \infty$, and $0 < \dim H_k < \infty$, according to the Fredholm alternative.

2. Let $u \in H_k$, $v \in H_l$ for $k \neq l$. Then $Su = \eta_k u$, $Sv = \eta_l v$ and so
$$\eta_k (u, v) = (Su, v) = (u, Sv) = \eta_l (u, v).$$
As $\eta_k \neq \eta_l$, we deduce $(u, v) = 0$. Consequently we see the subspaces $H_k$ and $H_l$ are orthogonal.

3. Now let $\tilde{H}$ be the smallest subspace of $H$ containing $H_0, H_1, \ldots$. Thus $\tilde{H} = \{\sum_{k=0}^m a_k u_k \mid m \in \{0, \ldots\}$, $u_k \in H_k$, $a_k \in \mathbb{R}\}$. We next demonstrate $\tilde{H}$ is dense in $H$. Clearly $S(\tilde{H}) \subset \tilde{H}$. Furthermore $S(\tilde{H}^\perp) \subset \tilde{H}^\perp$: indeed if $u \in \tilde{H}^\perp$ and $v \in \tilde{H}$, then $(Su, v) = (u, Sv) = 0$.

Now the operator $\tilde{S} = S|_{\tilde{H}^\perp}$ is compact and symmetric. In addition $\sigma(\tilde{S}) = \{0\}$, since any nonzero eigenvalue of $\tilde{S}$ would be an eigenvalue of $S$ as well. According to the lemma then, $(\tilde{S}u, u) = 0$ for all $u \in \tilde{H}^\perp$. But if $u, v \in \tilde{H}^\perp$,
$$2(\tilde{S}u, v) = (\tilde{S}(u + v), u + v) - (\tilde{S}u, u) - (\tilde{S}v, v) = 0.$$
Hence $\tilde{S} = 0$. Consequently $\tilde{H}^\perp \subset N(S) \subset \tilde{H}$, and so $\tilde{H}^\perp = \{0\}$. Thus $\tilde{H}$ is dense in $H$.

4. Choose an orthonormal basis for each subspace $H_k$ ($k = 0, \ldots$), noting that since $H$ is separable, $H_0$ has a countable orthonormal basis. We obtain thereby an orthonormal basis of eigenvectors. $\square$
\end{proof}

Most of these proofs are from \cite{Brezis1983}. See also \cite[Chapter 5]{GilbargTrudinger1983} and \cite{Yosida1980}.

\section{Measure Theory}
\label{sec:measure-theory}

This appendix provides a quick outline of some fundamentals of measure theory.

\subsection{Lebesgue Measure}
\label{sec:lebesgue-measure}

Lebesgue measure provides a way of describing the ``size'' or ``volume'' of certain subsets of $\mathbb{R}^n$.

\begin{definition}[$\sigma$-algebra]
\label{def:sigma-algebra}
\dcref{appE:E.1-def-1}
\group{Measure Theory}
A collection $\mathcal{M}$ of subsets of $\mathbb{R}^n$ is called a $\sigma$-algebra if

(i) $\emptyset, \mathbb{R}^n \in \mathcal{M}$,

(ii) $A \in \mathcal{M}$ implies $\mathbb{R}^n - A \in \mathcal{M}$,

and

(iii) if $\{A_k\}_{k=1}^{\infty} \subset \mathcal{M}$, then $\bigcup_{k=1}^{\infty} A_k, \bigcap_{k=1}^{\infty} A_k \in \mathcal{M}$.
\end{definition}

\begin{theorem}[Existence of Lebesgue measure and Lebesgue measurable sets]
\label{thm:existence-lebesgue-measure}
\dcref{appE:thm-1}
\group{Measure Theory}
\uses{def:sigma-algebra}
\textit{There exist a $\sigma$-algebra $\mathcal{M}$ of subsets of $\mathbb{R}^n$ and a mapping}
$$| \cdot | : \mathcal{M} \to [0, +\infty]$$
\textit{with the following properties:}

\begin{enumerate}
\item[(i)] \textit{Every open subset of $\mathbb{R}^n$, and thus every closed subset of $\mathbb{R}^n$, belong to $\mathcal{M}$.}
\item[(ii)] \textit{If $B$ is any ball in $\mathbb{R}^n$, then $|B|$ equals the $n$-dimensional volume of $B$.}
\item[(iii)] \textit{If $\{A_k\}_{k=1}^{\infty} \subset \mathcal{M}$ and the sets $\{A_k\}_{k=1}^{\infty}$ are pairwise disjoint, then}
$$(1) \quad \left| \bigcup_{k=1}^{\infty} A_k \right| = \sum_{k=1}^{\infty} |A_k| \quad \text{(``countable additivity")}.$$
\item[(iv)] \textit{If $A \subseteq B$, where $B \in \mathcal{M}$ and $|B| = 0$, then $A \in \mathcal{M}$ and $|A| = 0$.}
\end{enumerate}

\textbf{Notation.} The sets in $\mathcal{M}$ are called \textit{Lebesgue measurable sets} and $| \cdot |$ is $n$-dimensional \textit{Lebesgue measure}. $\square$
\end{theorem}

\begin{remark}[Basic properties of Lebesgue measure]
\label{rem:lebesgue-measure-properties}
\dcref{appE:E.1-rem-1}
\group{Measure Theory}
\uses{thm:existence-lebesgue-measure}
(i) From (ii) and (iii), we see that $|A|$ equals the volume of any set $A$ with piecewise smooth boundary.

(ii) We deduce from (1) that
$$(2) \quad |\emptyset| = 0,$$
and
$$(3) \quad \left| \bigcup_{k=1}^{\infty} A_k \right| \leq \sum_{k=1}^{\infty} |A_k| \quad \text{(``countable subadditivity")}$$
for any countable collection of measurable sets $\{A_k\}_{k=1}^{\infty}$.
\end{remark}

\begin{definition}[Almost everywhere]
\label{def:almost-everywhere}
\dcref{appE:E.1-def-2}
\group{Measure Theory}
\uses{thm:existence-lebesgue-measure}
If some property holds everywhere on $\mathbb{R}^n$, except for a measurable set with Lebesgue measure zero, we say the property holds \textit{almost everywhere}, abbreviated ``a.e.''. $\square$
\end{definition}

\subsection{Measurable Functions and Integration}
\label{sec:measurable-functions-integration}

\begin{definition}[Measurable function]
\label{def:measurable-function}
\dcref{appE:E.2-def-1}
\group{Measure Theory}
\uses{def:sigma-algebra}
Let $f : \mathbb{R}^n \to \mathbb{R}$. We say $f$ is a measurable function if
$$f^{-1}(U) \in \mathcal{M}$$
for each open subset $U \subset \mathbb{R}$.
\end{definition}

Note in particular that if $f$ is continuous, then $f$ is measurable. The sum and product of two measurable functions are measurable. In addition if $\{f_k\}_{k=0}^{\infty}$ are measurable functions, then so are $\limsup f_k$ and $\liminf f_k$.

\begin{theorem}[Egoroff's Theorem]
\label{thm:egoroffs-theorem}
\dcref{appE:thm-2}
\group{Measure Theory}
\uses{def:measurable-function,def:almost-everywhere}
Let $\{f_k\}_{k=1}^{\infty}$, $f$ be measurable functions, and
$$f_k \to f \quad \text{a.e. on } A,$$
where $A \subset \mathbb{R}^n$ is measurable, $|A| < \infty$. Then for each $\varepsilon > 0$ there exists a measurable subset $E \subset A$ such that
\begin{align*}
\text{(i)} & \quad |A - E| \le \varepsilon \\
\text{and} \\
\text{(ii)} & \quad f_k \to f \text{ uniformly on } E.
\end{align*}
\end{theorem}

Now if $f$ is a nonnegative, measurable function, it is possible, by an approximation of $f$ with simple functions, to define the Lebesgue integral
$$\int_{\mathbb{R}^n} f \, dx.$$
Cf. \S\ref{sec:banach-space-valued-functions} below. This agrees with the usual integral if $f$ is continuous or Riemann integrable. If $f$ is measurable, but is not necessarily nonnegative, we define
$$\int_{\mathbb{R}^n} f \, dx = \int_{\mathbb{R}^n} f^+ \, dx - \int_{\mathbb{R}^n} f^- \, dx,$$
provided at least one of the terms on the right hand side is finite. In this case we say $f$ is \textit{integrable}.

\begin{definition}[Summable function]
\label{def:summable-function}
\dcref{appE:E.2-def-2}
\group{Measure Theory}
\uses{def:measurable-function}
A measurable function $f$ is summable if
$$\int_{\mathbb{R}^n} |f| \, dx < \infty.$$
\end{definition}

\begin{remark}[Integrable versus summable terminology]
\label{rem:integrable-vs-summable-terminology}
\dcref{appE:E.2-rem-1}
\group{Measure Theory}
\uses{def:summable-function}
Note carefully our terminology: a measurable function is \textit{integrable} if it has an integral (which may equal $+\infty$ or $-\infty$) and is \textit{summable} if this integral is finite. $\square$
\end{remark}

\begin{definition}[Essential supremum]
\label{def:essential-supremum}
\dcref{appE:E.2-def-3}
\group{Measure Theory}
If the real-valued function $f$ is measurable, we define the \textit{essential supremum} of $f$ to be
$$\text{ess sup } f := \inf\{\mu \in \R \mid |\{f > \mu\}| = 0\}. \qquad \square$$
\end{definition}

\subsection{Convergence Theorems for Integrals}
\label{sec:convergence-theorems-for-integrals}

The Lebesgue theory of integration is especially useful since it provides the following powerful convergence theorems.

\begin{theorem}[Fatou's Lemma]
\label{thm:fatous-lemma}
\dcref{appE:thm-3}
\group{Measure Theory}
\uses{def:summable-function}
Assume the functions $\{f_k\}_{k=1}^{\infty}, f$ are nonnegative and summable. Then
$$\int_{\R^n} \liminf_{k \to \infty} f_k \, dx \leq \liminf_{k \to \infty} \int_{\R^n} f_k \, dx.$$
\end{theorem}

\begin{theorem}[Monotone Convergence Theorem]
\label{thm:monotone-convergence-theorem}
\dcref{appE:thm-4}
\group{Measure Theory}
\uses{def:summable-function}
Assume the functions $\{f_k\}_{k=1}^{\infty}$ are measurable, with
$$f_1 \leq f_2 \leq \cdots \leq f_k \leq f_{k+1} \leq \cdots.$$
Then
$$\int_{\R^n} \lim_{k \to \infty} f_k \, dx = \lim_{k \to \infty} \int_{\R^n} f_k \, dx.$$
\end{theorem}

\begin{theorem}[Dominated Convergence Theorem]
\label{thm:dominated-convergence-theorem}
\dcref{appE:thm-5}
\group{Measure Theory}
\uses{def:summable-function}
Assume the functions $\{f_k\}_{k=1}^{\infty}$ are integrable and
$$f_k \to f \quad \text{a.e.}$$
Suppose also
$$|f_k| \leq g \quad \text{a.e.},$$
for some summable function $g$. Then
$$\int_{\R^n} f_k \, dx \to \int_{\R^n} f \, dx.$$
\end{theorem}

\subsection{Differentiation}
\label{sec:differentiation}

An important fact is that a summable function is ``approximately continuous'' at almost every point.

\begin{theorem}[Lebesgue's Differentiation Theorem]
\label{thm:lebesgue-differentiation-theorem}
\dcref{appE:thm-6}
\group{Measure Theory}
\uses{def:summable-function,def:almost-everywhere}
Let $f : \mathbb{R}^n \to \mathbb{R}$ be \textit{locally summable}.

\begin{enumerate}
\item[(i)] Then for a.e. point $x_0 \in \mathbb{R}^n$,
$$\fint_{B(x_0,r)} f \, dx \to f(x_0) \quad \text{as } r \to 0.$$
\item[(ii)] In fact, for a.e. point $x_0 \in \mathbb{R}^n$,
$$(4) \quad \fint_{B(x_0,r)} |f(x) - f(x_0)| \, dx \to 0 \quad \text{as } r \to 0.$$
\end{enumerate}

A point $x_0$ at which (4) holds is called a \textit{Lebesgue point} of $f$.
\end{theorem}

\begin{remark}[$L^p$ version of the Lebesgue point theorem]
\label{rem:lebesgue-differentiation-lp}
\dcref{appE:E.4-rem-1}
\group{Measure Theory}
\uses{thm:lebesgue-differentiation-theorem}
More generally, if $f \in L^p_{\text{loc}}(\mathbb{R}^n)$ for some $1 \le p < \infty$, then for a.e. point $x_0 \in \mathbb{R}^n$ we have
$$\fint_{B(x_0,r)} |f(x) - f(x_0)|^p \, dx \to 0 \quad \text{as } r \to 0.$$
\end{remark}

\subsection{Banach Space-Valued Functions}
\label{sec:banach-space-valued-functions}

We extend the notions of measurability, integrability, etc. to mappings
$$f : [0, T] \to X$$
where $T > 0$ and $X$ is a real Banach space, with norm $\| \cdot \|$.

\begin{definition}[Simple function]
\label{def:simple-function-banach-valued}
\dcref{appE:E.5-def-1}
\group{Vector-Valued Integration}
\uses{def:sigma-algebra}
A function $s : [0, T] \to X$ is called \textit{simple} if it has the form
$$(5) \quad s(t) = \sum_{i=1}^{m} \chi_{E_i}(t) u_i \quad (0 \le t \le T),$$
where each $E_i$ is a \textit{Lebesgue measurable} subset of $[0, T]$ and $u_i \in X$ ($i = 1, \ldots, m$).
\end{definition}

\begin{definition}[Strongly measurable function]
\label{def:strongly-measurable-function}
\dcref{appE:E.5-def-2}
\group{Vector-Valued Integration}
\uses{def:simple-function-banach-valued,def:almost-everywhere}
A function $f : [0, T] \to X$ is \textit{strongly measurable} if there exist simple functions $s_k : [0, T] \to X$ such that
$$s_k(t) \to f(t) \quad \text{for a.e. } 0 \le t \le T.$$
\end{definition}

\begin{definition}[Weakly measurable function]
\label{def:weakly-measurable-function}
\dcref{appE:E.5-def-3}
\group{Vector-Valued Integration}
\uses{def:measurable-function,def:bounded-linear-functional-dual-space}
A function $f : [0, T] \to X$ is \textit{weakly measurable} if for each $u^* \in X^*$, the mapping $t \mapsto \langle u^*, f(t) \rangle$ is Lebesgue measurable.
\end{definition}

\begin{definition}[Almost separably valued]
\label{def:almost-separably-valued}
\dcref{appE:E.5-def-4}
\group{Vector-Valued Integration}
\uses{def:separable-space}
We say $f : [0,T] \to X$ is almost separably valued if there exists a subset $N \subset [0,T]$, with $|N| = 0$, such that the set $\{f(t) \mid t \in [0,T] - N\}$ is separable.
\end{definition}

\begin{theorem}[Pettis]
\label{thm:pettis-theorem}
\dcref{appE:thm-7}
\group{Vector-Valued Integration}
\uses{def:strongly-measurable-function,def:weakly-measurable-function,def:almost-separably-valued}
The mapping $f : [0,T] \to X$ is strongly measurable if and only if $f$ is weakly measurable and is almost separably valued.
\end{theorem}

\begin{definition}[Integral of a simple function]
\label{def:integral-of-simple-function}
\dcref{appE:E.5-def-5}
\group{Vector-Valued Integration}
\uses{def:simple-function-banach-valued}
If $s(t) = \sum_{i=1}^m \chi_{E_i}(t)u_i$ is simple, we define
$$(6) \quad \int_0^T s(t) \, dt := \sum_{i=1}^m |E_i|u_i.$$
\end{definition}

\begin{definition}[Summable Banach space-valued function]
\label{def:summable-banach-valued-function}
\dcref{appE:E.5-def-6}
\group{Vector-Valued Integration}
\uses{def:integral-of-simple-function}
We say $f : [0,T] \to X$ is summable if there exists a sequence $\{s_k\}_{k=1}^\infty$ of simple functions such that
$$(7) \quad \int_0^T \|s_k(t) - f(t)\| \, dt \to 0 \text{ as } k \to \infty.$$
\end{definition}

\begin{definition}[Bochner integral]
\label{def:bochner-integral}
\dcref{appE:E.5-def-7}
\group{Vector-Valued Integration}
\uses{def:summable-banach-valued-function,def:integral-of-simple-function}
If $f$ is summable, we define
$$(8) \quad \int_0^T f(t) \, dt = \lim_{k \to \infty} \int_0^T s_k(t) \, dt.$$
\end{definition}

\begin{theorem}[Bochner]
\label{thm:bochner-theorem}
\dcref{appE:thm-8}
\group{Vector-Valued Integration}
\uses{def:strongly-measurable-function,def:summable-banach-valued-function,def:bochner-integral,def:summable-function}
A strongly measurable function $f : [0,T] \to X$ is summable if and only if $t \mapsto \|f(t)\|$ is summable. In this case
$$\left\| \int_0^T f(t) \, dt \right\| \leq \int_0^T \|f(t)\| \, dt,$$
and
$$\left\langle u^*, \int_0^T f(t) \, dt \right\rangle = \int_0^T \langle u^*, f(t) \rangle \, dt$$
for each $u^* \in X^*$.
\end{theorem}

A good book for measure theory is \cite{Folland1984}. See \cite[Chapter V, Sections 4--5]{Yosida1980} for proofs of \cref{thm:pettis-theorem} and \cref{thm:bochner-theorem}.
